\section{从关门到京师}

顺治元年四月，山海关不是一扇一经开启便可直达天下的门，而是一处把几种危机压在一起的关隘。李自成入北京之后，关宁军同大顺之间的兵饷、统属与互信已经破裂；多尔衮则把这一裂缝变成八旗越过关隘的机会。清军与吴三桂部在山海关击败大顺军，因而取得进入华北的通道。这一结果可以立为编年事实；但若把后来流行的“乞师”故事逐句当作战前文件，把各军人数、会师次序和战场上的每一次转折都写成定论，证据便承受不住这样的重压。

《清世祖实录》中的军报把胜利纳入摄政政权的行动叙述，清初内国史院保存的满文关宁军务档案则让人看见命令、呈报和接收如何在文书链中移动。二者首先是清廷留下的行政记录：它们能证明朝廷收到什么消息、以何种语句处分军政，却不能自动代替关内各军的全部经验。后出的《清史稿》适合用来互校官名和大事顺序，但它是民国初年未定稿的官修史书，不是顺治四月的现场记录。吴三桂求援文书的原貌、各方兵额以及战场过程，仍须保留为由清廷叙事和后出文本共同塑形的问题。

五月入北京，也不应被压缩成“占领京城”四字。大顺撤离、明廷解体之后，城中留下的是仓储、官署、粮饷、逃散的人群和等待选择的官员。多尔衮率军入城，接收的是明朝京师的行政与礼制空间；到冬季，摄政王又以安民、征发、招纳明官和整顿官署的命令，尝试把军队占有的城池变成可以传递文书的治所。诏令的发布是可靠的，题本中可见的交接也证明部分官署重新运转；可是受任者是否到任、一纸命令是否越过城门和驿路、地方是否照办，不能从诏书的语气倒推出去。

北京地方记事和《甲申传信录》保存了另一种尺度：它们所见的是城内消息、秩序与恐惧，而非摄政政权所要展示的承接。它们同样有作者位置、传抄与见闻范围的限制，却正可防止以中央档案的沉默替代城市生活。顺治帝九月自盛京抵达北京、举行入城和定都仪式，标志着皇帝驻跸与部院中心向北京集中；盛京却仍保有满洲政治和象征上的重量。仪式建立的是持续的中央名义，并不等于各地已经承认，尤其不等于南方战场已经结束。

\section{摄政并非一条畅通的命令线}

多尔衮的摄政，是清初征服得以调动兵马、官员和名义资源的枢纽，但不是一架无阻运转的机器。京师需要粮饷，部院需要熟悉旧制的书吏和降官，前线又不断索要兵丁、漕运和军需。清廷档案最擅长呈现这条由奏报、批示、部议到任命的路径；地方档案、方志和南方对手的材料则不断提醒人们，路径在何处中断、被拖延，或被战事改写。把“招降”写成一个完成状态，会遮蔽降官仍须在家族、旧部、军饷和地方安全之间反复选择的处境。

顺治七年十二月，多尔衮在古北口外行猎途中去世。其死期与丧仪可以由《清世祖实录》和内国史院的遗命、丧仪档相互钉住；死因以及死后政治清算的动机，却不能只照后来定性的语言复述。次年顺治帝亲政，摄政安排撤除，皇帝周边的任官与决策网络重新组合。然而，署在皇帝名下的命令仍要依靠宗室、旗务、满汉官僚与南方将领落实。摄政结束，不是中央能力忽然圆满的时刻。

这一点在顺治帝去世后更为清楚。顺治十八年，八岁的玄烨即位，索尼、苏克萨哈、遏必隆、鳌拜辅政。继位与四辅政大臣的名单见于《清世祖实录》、起居注和遗诏档；至康熙四年鳌拜影响上升、康熙六年索尼去世、康熙八年鳌拜被拘捕，则应分别阅读《清圣祖实录》、起居注与内阁案卷。尤其不能把1669年的结果倒灌回此前：后来的鳌拜案和人物传记最容易把本来由议政、旗务、人事和军费共同构成的辅政政治，写成一人独揽权力的必然前史。康熙拘捕鳌拜使皇帝能够更直接地统摄人事军政，却也仍须在既有制度与旗人贵族网络中行使这种权力。

\section{南京失守以后：城池、朝廷与身体}

清军南下的速度使京师接收和江南战争看似连成一线，实际却是一系列由水陆节点、粮道和地方选择构成的战区。顺治二年春夏之交，清军沿运河、淮扬推进，攻取南京；弘光帝朱由崧不久在芜湖一带被俘。清方的《清世祖实录》与《清史稿·世祖本纪》可以确立城陷、押解这一主干，南明弘光朝奏疏、《南明史料》所收文书和江宁地方记事则保留了另一面：朝廷如何失去护卫与指挥，官员和军队怎样逃散、观望或改易归属。前者是胜利政权收到的战报，后者多是败亡朝廷及城市中的急迫文字；两类材料不能互相注销，也不应被合成为一种全知叙述。

南京失去首都地位，并没有令“明亡”在南方成为一个同时发生的事件。弘光名义瓦解后，鲁王集团在浙东和海上活动，隆武朝廷在福建建立，绍武与永历又在两广相继争夺名义中心。清军取得杭州、绍兴等城，或攻入广州，并不自动说明一省的道路、岛屿、山地和县城都已归于同一控制。清廷军报常以府城、俘获和招降排列战果；《福建通志》《广东通志》及各地府县志可以补足部分地名和后续秩序，却大多成书较晚、也会沿袭胜败叙事。南明文集和《永历实录》把流动朝廷、将领联络与退守路线连成另一条线索，但它们也服务于保存正统、追责或追忆的需要。史实的难处不在于挑选哪一方“说真话”，而在于说明每种文本首先能够证明什么。

顺治二年六月重申的薙发易服命令，使政治归属进入身体的可见处。实录中的命令证明清廷以发式服制辨认新旧政治身份；兵科题本和江南府县志的相关条目，则可见命令在若干地方与征发、守城和治安纠缠在一起。它不能证明全国在同一日、以同一种方式服从。有人遵从，有人逃亡，有人把发式当作拒绝交付城池或粮饷的标识；不同府县的速度、强度与后果必须逐地追问。

因此，扬州、江阴和1650年的广州特别不宜用一个整齐的死亡数字盖棺。扬州失陷、江阴围城失守、广州再陷，以及攻城后发生严重杀戮、掳掠和损毁，都是必须写入征服史的事实。可是清军战报所陈的战果，不是伤亡统计；《扬州十日记》《江阴城守纪》、广州地方笔记和后出的府县志，保存了幸存者、编纂者或地方共同体的恐惧与纪念，也各有成书、传抄、修辞和数字来源的问题。“十日”、围城天数和死亡人数不应被伪装成精确普查，更不能由一部文本推算江南或广东整体的死亡。较为诚实的写法是同时承认暴力及其长期记忆，说明现有材料不能可靠合算总数，并把不同数字标回提出它的作者、地点与文类。

\section{南方并未在一役中归入清廷}

1646年隆武帝在闽北被俘杀、绍武朝在广州结束之后，朱由榔在肇庆即位，永历年号成为两广和西南多支力量可以借用的名义。清廷军报将这看作继续清剿的对象；《永历实录》、永历朝奏疏和地方志所呈现的，却是一个不断移动、依赖地方军队、饷源与山海通道的朝廷。1648年金声桓在南昌、李成栋在广东改奉南明，说明此前的“接收”未能冻结将领与地方的选择。1649年南昌被攻取、1650年广州再陷固然削弱了这些反转的中心，却不能把战后的税源、流民和港口恢复一并写成已经解决。

永历六年，李定国攻取桂林、在衡州击败清军，迫使清廷重新估量湘桂道路、军需和机动兵力。双方的兵数、死伤和将领死因在清廷军报与南明记录之间多有不同；可以肯定战场结果，不能替任何一方把报捷数字升级为无争议的计数。此后孙可望与永历朝的冲突及其降清，削弱了南明的军事协作；清军自贵州、四川等方向入滇，永历朝廷离开昆明，终至进入缅甸。这里又有一层不能省略的证据边界：清方材料擅长记录追捕与军政调度，缅甸宫廷的《琉璃宫史》传统和欧洲旅行记所保留的接待、路线与交涉，并不必然与之重合。把缅甸只写作清军追击地图上的空白，既抹掉了东吁王朝的选择，也误解了流亡朝廷的补给条件。

1659年郑成功北上围攻南京而后撤退，显示长江下游并非已无战事；翌年清军经营云南城镇和驿路，也说明占有昆明不等于控制边地。1662年，朱由榔由缅甸被押回后处死，南明皇帝的名义中心至此终结。中缅材料对于移交过程、处死地点和责任分配仍须并读。夔东抗清据点直到1664年才被清军攻破，郑氏又在海上和台湾维持政权；这些后续并非枝节，而是对“1644年入关即全国统一”这一方便说法的校正。

从山海关到北京，从南京到西南边地，清初形成的不是一条单向延伸的疆界，而是军队、文书、仪式、降附、逃亡与再选择彼此牵动的过程。清廷档案使我们能够看见命令如何被制作和宣告；官修史书为后来者编排出帝国的时间；地方材料留下城池和社区承受的断片；南明叙事则保存失败者仍在行动的政治地理。把它们放在同一页上，不是为了把分歧磨平，而是为了让征服、摄政与整合都回到它们当时尚未完成的样子。


\subsection{事件链与材料限度}
关隘、京城与官署的材料并不构成一条无摩擦的接收线。以下事件保留各自的行动者、空间、后果与证据限度。

\paragraph{顺治元年五月（1644年6月前后）：多尔衮率军进入北京并接收明朝京城的行政与礼制空间}
\eventanchor{EQ-1644-BEIJING-ENTRY}{多尔衮率军进入北京并接收明朝京城的行政与礼制空间}。本节点叙述该事件本身。清摄政政权与北京在京师与多地军政线路相接的尺度的处境首先受大顺撤离与明廷崩溃留下军政真空，清军需获取京师粮饷、官僚和象征中心制约。《清世祖实录》顺治元年五月；《清史稿·世祖本纪》；《甲申传信录》及北京地方记事留下可核的线索；可辨的后果是清廷以北京为治所，开始将军事占领转化为对明朝制度资源的继承和重组。原有置信注记：入城与接收可靠；城内秩序、迎降者范围和掠夺状况须以不同见证互校

\paragraph{顺治元年五月（1644年6月前后）：多尔衮率军进入北京并接收明朝京城的行政与礼制空间}
\eventanchor{EQ-1644-BEIJING-ENTRY-AFTERMATH}{多尔衮率军进入北京并接收明朝京城的行政与礼制空间} 置于京师与多地军政线路相接的尺度中阅读，本节点追踪「多尔衮率军进入北京并接收明朝京城的行政与礼制空间」之后可见的余波。大顺撤离与明廷崩溃留下军政真空，清军需获取京师粮饷、官僚和象征中心构成行动者清摄政政权与北京必须面对的条件。取证从《清世祖实录》顺治元年五月；《清史稿·世祖本纪》；《甲申传信录》及北京地方记事出发，因而只能把变化谨慎地连到清廷以北京为治所，开始将军事占领转化为对明朝制度资源的继承和重组。原有置信注记：入城与接收可靠；城内秩序、迎降者范围和掠夺状况须以不同见证互校

\paragraph{顺治元年五月（1644年6月前后）：多尔衮率军进入北京并接收明朝京城的行政与礼制空间}
\eventanchor{EQ-1644-BEIJING-ENTRY-DECISION}{多尔衮率军进入北京并接收明朝京城的行政与礼制空间} 标记本节点定位围绕「多尔衮率军进入北京并接收明朝京城的行政与礼制空间」形成的处置与调度记录，涉及清摄政政权与北京。京师与多地军政线路相接的尺度不是抽象背景：大顺撤离与明廷崩溃留下军政真空，清军需获取京师粮饷、官僚和象征中心。《清世祖实录》顺治元年五月；《清史稿·世祖本纪》；《甲申传信录》及北京地方记事所见的顺序随后指向清廷以北京为治所，开始将军事占领转化为对明朝制度资源的继承和重组，但原有置信注记：入城与接收可靠；城内秩序、迎降者范围和掠夺状况须以不同见证互校

\paragraph{顺治元年五月（1644年6月前后）：多尔衮率军进入北京并接收明朝京城的行政与礼制空间}
在京师与多地军政线路相接的尺度，\eventanchor{EQ-1644-BEIJING-ENTRY-PRELUDE}{多尔衮率军进入北京并接收明朝京城的行政与礼制空间} 让清摄政政权与北京的一个选择或压力有了可查的坐标。本节点回看「多尔衮率军进入北京并接收明朝京城的行政与礼制空间」之前已可见的条件。前因可写为大顺撤离与明廷崩溃留下军政真空，清军需获取京师粮饷、官僚和象征中心；《清世祖实录》顺治元年五月；《清史稿·世祖本纪》；《甲申传信录》及北京地方记事限定了记录，后续变化是清廷以北京为治所，开始将军事占领转化为对明朝制度资源的继承和重组。原有置信注记：入城与接收可靠；城内秩序、迎降者范围和掠夺状况须以不同见证互校

\paragraph{顺治元年冬（1644年）：多尔衮以摄政王名义整顿京师官署、招降明官并发布安民与征发命令}
\eventanchor{EQ-1644-REGENT-ORDERS}{多尔衮以摄政王名义整顿京师官署、招降明官并发布安民与征发命令} 所关联的是清摄政政权在京师与多地军政线路相接的尺度中的行动。本节点叙述该事件本身。它从新政权需要恢复行政传递、维持京师粮饷并吸纳明代文官技术展开，借《清世祖实录》顺治元年十月至十二月；《清史稿·职官志》；中国第一历史档案馆藏内阁大库顺治元年题本〔京师官署交接〕进入可检验的年序，并以中央可用的部院和任官网络开始恢复，但忠诚、财力与军纪继续限制覆盖面影响下一段安排。原有置信注记：命令发布可靠；接受任官和地方执行不能从诏令直接推出

\paragraph{顺治元年冬（1644年）：多尔衮以摄政王名义整顿京师官署、招降明官并发布安民与征发命令}
京师与多地军政线路相接的尺度中的\eventanchor{EQ-1644-REGENT-ORDERS-AFTERMATH}{多尔衮以摄政王名义整顿京师官署、招降明官并发布安民与征发命令} 不能离开清摄政政权来解释。本节点追踪「多尔衮以摄政王名义整顿京师官署、招降明官并发布安民与征发命令」之后可见的余波。新政权需要恢复行政传递、维持京师粮饷并吸纳明代文官技术说明其形成的条件；《清世祖实录》顺治元年十月至十二月；《清史稿·职官志》；中国第一历史档案馆藏内阁大库顺治元年题本〔京师官署交接〕提供来源路径，中央可用的部院和任官网络开始恢复，但忠诚、财力与军纪继续限制覆盖面则是材料能够支持的近时结果。原有置信注记：命令发布可靠；接受任官和地方执行不能从诏令直接推出

\paragraph{顺治元年冬（1644年）：多尔衮以摄政王名义整顿京师官署、招降明官并发布安民与征发命令}
\eventanchor{EQ-1644-REGENT-ORDERS-DECISION}{多尔衮以摄政王名义整顿京师官署、招降明官并发布安民与征发命令} 把清摄政政权放回京师与多地军政线路相接的尺度的道路、城镇、边界或制度网络。本节点定位围绕「多尔衮以摄政王名义整顿京师官署、招降明官并发布安民与征发命令」形成的处置与调度记录。新政权需要恢复行政传递、维持京师粮饷并吸纳明代文官技术。本段采用《清世祖实录》顺治元年十月至十二月；《清史稿·职官志》；中国第一历史档案馆藏内阁大库顺治元年题本〔京师官署交接〕核对其顺序，随后可追到中央可用的部院和任官网络开始恢复，但忠诚、财力与军纪继续限制覆盖面。原有置信注记：命令发布可靠；接受任官和地方执行不能从诏令直接推出

\paragraph{顺治元年冬（1644年）：多尔衮以摄政王名义整顿京师官署、招降明官并发布安民与征发命令}
本节点回看「多尔衮以摄政王名义整顿京师官署、招降明官并发布安民与征发命令」之前已可见的条件，故\eventanchor{EQ-1644-REGENT-ORDERS-PRELUDE}{多尔衮以摄政王名义整顿京师官署、招降明官并发布安民与征发命令} 不只是一个年份标签。清摄政政权在京师与多地军政线路相接的尺度承受的压力来自新政权需要恢复行政传递、维持京师粮饷并吸纳明代文官技术。《清世祖实录》顺治元年十月至十二月；《清史稿·职官志》；中国第一历史档案馆藏内阁大库顺治元年题本〔京师官署交接〕可回查该节点；中央可用的部院和任官网络开始恢复，但忠诚、财力与军纪继续限制覆盖面是其进入后续年序的方式。原有置信注记：命令发布可靠；接受任官和地方执行不能从诏令直接推出

\paragraph{顺治元年四月（1644年5月，换算待核）：清军与吴三桂在山海关击败李自成军，打开进入华北的军事通道}
\eventanchor{EQ-1644-SHANHAIGUAN-AFTERMATH}{清军与吴三桂在山海关击败李自成军，打开进入华北的军事通道} 所见的行动者为清摄政政权、吴三桂与大顺，空间尺度为京师与多地军政线路相接的尺度。本节点追踪「清军与吴三桂在山海关击败李自成军，打开进入华北的军事通道」之后可见的余波。李自成入北京后关宁军与大顺在兵饷、统属和互信上破裂，多尔衮把危机转为入关机会使它成为具体事件链的一环；《清世祖实录》顺治元年四月山海关军报；《清史稿·世祖本纪》；《清初内国史院满文档案》顺治元年四月关宁军务档与清军获得穿越关隘的行动空间，北京政权争夺转为华北与南方的多战区征服。原有置信注记：核心战事可靠；吴三桂求援文书原貌、各军兵数和战场过程受清廷与后出叙事塑形

\paragraph{顺治元年四月（1644年5月，换算待核）：清军与吴三桂在山海关击败李自成军，打开进入华北的军事通道}
从京师与多地军政线路相接的尺度看，\eventanchor{EQ-1644-SHANHAIGUAN-DECISION}{清军与吴三桂在山海关击败李自成军，打开进入华北的军事通道} 留下本节点定位围绕「清军与吴三桂在山海关击败李自成军，打开进入华北的军事通道」形成的处置与调度记录的材料坐标。清摄政政权、吴三桂与大顺所面对的条件是李自成入北京后关宁军与大顺在兵饷、统属和互信上破裂，多尔衮把危机转为入关机会。《清世祖实录》顺治元年四月山海关军报；《清史稿·世祖本纪》；《清初内国史院满文档案》顺治元年四月关宁军务档支持的结果为清军获得穿越关隘的行动空间，北京政权争夺转为华北与南方的多战区征服。原有置信注记：核心战事可靠；吴三桂求援文书原貌、各军兵数和战场过程受清廷与后出叙事塑形

\paragraph{顺治元年四月（1644年5月，换算待核）：清军与吴三桂在山海关击败李自成军，打开进入华北的军事通道}
\eventanchor{EQ-1644-SHANHAIGUAN-PRELUDE}{清军与吴三桂在山海关击败李自成军，打开进入华北的军事通道} 记录清摄政政权、吴三桂与大顺在京师与多地军政线路相接的尺度的一次变化。本节点回看「清军与吴三桂在山海关击败李自成军，打开进入华北的军事通道」之前已可见的条件。其发生与李自成入北京后关宁军与大顺在兵饷、统属和互信上破裂，多尔衮把危机转为入关机会相连；《清世祖实录》顺治元年四月山海关军报；《清史稿·世祖本纪》；《清初内国史院满文档案》顺治元年四月关宁军务档把事件系回来源，清军获得穿越关隘的行动空间，北京政权争夺转为华北与南方的多战区征服显示压力如何向后转移。原有置信注记：核心战事可靠；吴三桂求援文书原貌、各军兵数和战场过程受清廷与后出叙事塑形

\paragraph{顺治元年九月（1644年10月前后）：顺治帝自盛京抵北京，清廷举行定都与皇帝入城仪式}
\eventanchor{EQ-1644-SHUNZHI-BEIJING}{顺治帝自盛京抵北京，清廷举行定都与皇帝入城仪式} 的地点与行动范围属于京师与多地军政线路相接的尺度，行动者是清。本节点叙述该事件本身。摄政政权虽占北京，仍需皇帝亲临与礼制重置来建立持续中央名义。读《清世祖实录》顺治元年九月；《清史稿·世祖本纪》；《清初内国史院满文档案》北京迁都礼仪档时可见其记录边界，最小后果只能写作北京成为皇帝驻跸与部院中心，盛京仍保有满洲政治和象征地位。原有置信注记：日期与仪式主干可靠；官修礼仪记录不能等同各地承认

\paragraph{顺治元年九月（1644年10月前后）：顺治帝自盛京抵北京，清廷举行定都与皇帝入城仪式}
\eventanchor{EQ-1644-SHUNZHI-BEIJING-AFTERMATH}{顺治帝自盛京抵北京，清廷举行定都与皇帝入城仪式} 是清在京师与多地军政线路相接的尺度留下的编年标记。本节点追踪「顺治帝自盛京抵北京，清廷举行定都与皇帝入城仪式」之后可见的余波。它从摄政政权虽占北京，仍需皇帝亲临与礼制重置来建立持续中央名义走来；《清世祖实录》顺治元年九月；《清史稿·世祖本纪》；《清初内国史院满文档案》北京迁都礼仪档固定可证部分，北京成为皇帝驻跸与部院中心，盛京仍保有满洲政治和象征地位构成下一步的可见结果。原有置信注记：日期与仪式主干可靠；官修礼仪记录不能等同各地承认

\paragraph{顺治元年九月（1644年10月前后）：顺治帝自盛京抵北京，清廷举行定都与皇帝入城仪式}
本节点定位围绕「顺治帝自盛京抵北京，清廷举行定都与皇帝入城仪式」形成的处置与调度记录：\eventanchor{EQ-1644-SHUNZHI-BEIJING-DECISION}{顺治帝自盛京抵北京，清廷举行定都与皇帝入城仪式}。清在京师与多地军政线路相接的尺度的处境可由摄政政权虽占北京，仍需皇帝亲临与礼制重置来建立持续中央名义说明。《清世祖实录》顺治元年九月；《清史稿·世祖本纪》；《清初内国史院满文档案》北京迁都礼仪档是本段材料入口，因而只能据此追踪到北京成为皇帝驻跸与部院中心，盛京仍保有满洲政治和象征地位。原有置信注记：日期与仪式主干可靠；官修礼仪记录不能等同各地承认

\paragraph{顺治元年九月（1644年10月前后）：顺治帝自盛京抵北京，清廷举行定都与皇帝入城仪式}
\eventanchor{EQ-1644-SHUNZHI-BEIJING-PRELUDE}{顺治帝自盛京抵北京，清廷举行定都与皇帝入城仪式} 发生在京师与多地军政线路相接的尺度，牵动清。本节点回看「顺治帝自盛京抵北京，清廷举行定都与皇帝入城仪式」之前已可见的条件。摄政政权虽占北京，仍需皇帝亲临与礼制重置来建立持续中央名义是其结构性条件；《清世祖实录》顺治元年九月；《清史稿·世祖本纪》；《清初内国史院满文档案》北京迁都礼仪档留下的证据把读者带到北京成为皇帝驻跸与部院中心，盛京仍保有满洲政治和象征地位。原有置信注记：日期与仪式主干可靠；官修礼仪记录不能等同各地承认